% ============================================================
%  Chapter 1 – Introduction
% ============================================================
\chapter{Introduction}

\section{Overview}

The coastal belt of Kerala, particularly the Kanayannur district, is highly vulnerable to a meteorological phenomenon locally known as \textit{Kallakkadal} — a long-period ocean swell surge that reaches the shore without perceptible warning, causing severe casualties, damage to fishing boats, and coastal property destruction.
Unlike tsunamis, Kallakkadal events are driven by distant storm systems in the Arabian Sea or Indian Ocean, generating long-wavelength swells that travel thousands of kilometres before amplifying near shore due to bathymetric shoaling.

Traditional early-warning mechanisms in Kerala rely on broadcast radio announcements or informal tip networks among fishing communities, offering little quantitative lead time and no real-time confirmation.
WaveGuard V4 proposes a technology-driven, sensor-fused, and publicly accessible early-warning system that addresses these gaps through an integrated stack of Internet-of-Things hardware and modern web technologies.

\section{Problem Statement}

Coastal communities engaged in fishing and maritime commerce operate at daily risk during the Arabian Sea swell season (May–November).
The core challenges are:

\begin{enumerate}[label=\roman*.]
  \item \textbf{Absence of real-time quantitative monitoring:} No publicly available near-shore buoy network provides live wave data for the Kanayannur coastal zone.
  \item \textbf{High false-alarm rate:} Existing informal warning systems often generate false alarms from boat wakes or other transient disturbances, leading to \textit{alarm fatigue} among residents.
  \item \textbf{Language accessibility:} Warning dissemination in English alone excludes non-literate Malayalam-speaking fishing households.
  \item \textbf{Alert latency:} Manual relay of warnings introduces unacceptable delays in life-safety scenarios.
\end{enumerate}

\section{Objectives}

The primary objectives of the WaveGuard V4 project are:

\begin{enumerate}[label=\arabic*.]
  \item Design and deploy a low-cost ESP32-based marine buoy instrumented with an MPU-6050 six-axis inertial measurement unit.
  \item Develop a FastAPI backend service implementing a multi-factor data-fusion algorithm combining buoy accelerometer readings with satellite-derived significant wave height from the Open-Meteo API.
  \item Implement a three-tier alert classification system (NORMAL, WATCH, WARNING) with defined numerical thresholds to minimise false positives.
  \item Build a real-time React + Vite web frontend that delivers live telemetry, alert status, and historical trend charts to the public.
  \item Provide a secure JWT-authenticated admin terminal for system management and status monitoring.
  \item Support bilingual output in English and Malayalam to serve local fishing communities.
  \item Lay the architectural groundwork for scaling to a multi-buoy mesh network.
\end{enumerate}

\section{Scope of the Project}

WaveGuard V4 operates within the following scope:

\begin{itemize}
  \item \textbf{Geographic scope:} Primary deployment zone is the Kanayannur coastal segment of Ernakulam district, Kerala (approximately 9.9–10.0°N, 76.2–76.3°E).
  \item \textbf{Sensing scope:} A single prototype buoy measuring tri-axial acceleration at up to 100 Hz and satellite wave height updated every ten minutes.
  \item \textbf{User scope:} Public safety dashboard accessible by any internet-connected device; admin terminal restricted to authorised personnel.
  \item \textbf{Current limitations:} SMS and push notifications are architecturally planned but not yet operationally deployed; multi-buoy networking is deferred to future work.
\end{itemize}

\section{Sustainable Development Goal Mapping}

WaveGuard V4 directly contributes to three United Nations Sustainable Development Goals \cite{undp_sdg}:

\begin{description}
  \item[\textbf{SDG 11 – Sustainable Cities and Communities}] The system provides early-warning infrastructure that strengthens the resilience and adaptive capacity of coastal communities to natural hazards, directly supporting SDG Target 11.5 (reduce disaster-related mortality and economic loss).
  \item[\textbf{SDG 13 – Climate Action}] By integrating satellite climate data with in-situ sensing, WaveGuard V4 builds community capacity to understand and respond to climate-driven sea-state anomalies (SDG Target 13.1).
  \item[\textbf{SDG 14 – Life Below Water}] The project supports safe maritime activity in coastal waters, reduces the risk to fisherfolk, and indirectly promotes sustainable ocean resource utilisation by enabling data-driven fishing decisions (SDG Target 14.2).
\end{description}

\section{Organisation of the Report}

The remainder of this report is organised as follows.
Chapter~\ref{chap:litsurvey} reviews related work on coastal monitoring and IoT-based early-warning systems.
Chapter~\ref{chap:background} describes the background technologies employed.
Chapter~\ref{chap:reqspec} presents the functional and non-functional requirements.
Chapter~\ref{chap:design} details the system design and implementation.
Chapter~\ref{chap:perfeval} presents the performance evaluation methodology and results.
Chapter~\ref{chap:results} presents screenshots and key metrics.
Chapter~\ref{chap:future} discusses future enhancements.
Chapter~\ref{chap:conclusion} concludes the report.
